\chapter{The front end: lexer, parser, semantic analysis}
\label{ch:frontend}

\section{Lexical analysis}
\subsection{Plain-language view}
The lexer turns the character stream \verb|let x: int = 42;| into the token stream \code{let}, \code{x}, \code{:}, \code{int}, \code{=}, \code{42}, \code{;}. It also discards whitespace and comments and remembers where each token started, so that later error messages can point at the exact line and column.

\subsection{Token classes}
\begin{table}[H]\centering\small
\begin{tabular}{@{}lp{0.62\textwidth}@{}}\toprule
\textbf{Class} & \textbf{Pattern}\\ \midrule
Identifier & \verb|[A-Za-z_][A-Za-z0-9_]*| (unless a keyword)\\
Keywords & \code{fn let if else while return print int float bool true false string new}\\
Integer literal & \verb|[0-9]+| (must fit in a signed 64-bit integer)\\
Float literal & \verb|[0-9]+ . [0-9]+| (digits are required on both sides of the dot)\\
String literal & \verb|"..."| with escapes \verb|\n \t \\ \"|; must close on the same line\\
Operators & \verb|+ - * / % < <= > >= == != && || ! = ->|\\
Punctuation & \verb|( ) { } [ ] , ; :|\\
Comments & \verb|// ...| to end of line; \verb|/* ... */| (not nested)\\ \bottomrule
\end{tabular}
\caption{Lexical grammar.}
\end{table}

The scanner is a hand-written equivalent of a deterministic finite automaton: it looks at one character of look-ahead and applies the longest-match rule, which is why \verb|<=| is a single token. For the program \verb|fn main() { print(1 + 2 * 3); }| it produces the stream in Table~\ref{tab:tokens} (first 20 of \SampleTokens{} tokens shown; the \emph{terminal} column is the name the parser's grammar uses).

\begin{table}[H]\centering\small
\begin{tabular}{@{}rllrl@{}}\toprule
\textbf{\#} & \textbf{Category} & \textbf{Lexeme} & \textbf{Line:Col} & \textbf{Grammar terminal}\\ \midrule
0 & KEYWORD & \texttt{fn} & 1:1 & \texttt{fn}\\
1 & IDENTIFIER & \texttt{main} & 1:4 & \texttt{ID}\\
2 & PUNCTUATION & \texttt{(} & 1:8 & \texttt{(}\\
3 & PUNCTUATION & \texttt{)} & 1:9 & \texttt{)}\\
4 & PUNCTUATION & \texttt{\{} & 1:11 & \texttt{\{}\\
5 & KEYWORD & \texttt{print} & 1:13 & \texttt{print}\\
6 & PUNCTUATION & \texttt{(} & 1:18 & \texttt{(}\\
7 & INT\_LITERAL & \texttt{1} & 1:19 & \texttt{INT\_LIT}\\
8 & OPERATOR & \texttt{+} & 1:21 & \texttt{+}\\
9 & INT\_LITERAL & \texttt{2} & 1:23 & \texttt{INT\_LIT}\\
10 & OPERATOR & \texttt{*} & 1:25 & \texttt{*}\\
11 & INT\_LITERAL & \texttt{3} & 1:27 & \texttt{INT\_LIT}\\
12 & PUNCTUATION & \texttt{)} & 1:28 & \texttt{)}\\
13 & PUNCTUATION & \texttt{;} & 1:29 & \texttt{;}\\
14 & PUNCTUATION & \texttt{\}} & 1:31 & \texttt{\}}\\
15 & EOF & \texttt{} & 1:32 & \texttt{\$}\\
\bottomrule
\end{tabular}

\caption{Real token stream for \texttt{fn main() \{ print(1 + 2 * 3); \}}.}\label{tab:tokens}
\end{table}

\subsection{Lexical errors}
Eight error codes cover malformed input; each has a position range, an explanation and a suggested fix, e.g. \texttt{E-LEX-006} (``Integer literal does not fit in 64 bits --- use a float literal'') or \texttt{E-LEX-007} (unterminated string). The scanner always makes progress after an error, so one run reports every lexical problem in the file.

\section{Syntax analysis: the canonical LR(1) parser}
\subsection{The grammar}
The grammar (Listing~\ref{lst:grammar}) has \LRProductions{} productions over \LRTerminals{} terminals and \LRNonterminals{} nonterminals. Three design decisions keep it unambiguous and conflict-free:

\begin{itemize}
  \item \textbf{Stratified expressions.} Precedence is encoded by one nonterminal per level (\code{OrExpr}, \code{AndExpr}, \code{EqExpr}, \code{RelExpr}, \code{AddExpr}, \code{MulExpr}, \code{UnaryExpr}, \code{Postfix}, \code{Primary}); every binary level is left-recursive, which gives left associativity.
  \item \textbf{Blocks are mandatory} after \code{if}/\code{while}, which removes the dangling-else ambiguity.
  \item \textbf{Types are structured} as \code{Type -> ElemType | ElemType [ ] | string} so that \code{int[]} parses without ambiguity against indexing.
\end{itemize}

\subsection{Building the tables}
The generator (\code{compiler/src/parser/lr1.cpp}) implements exactly the algorithm of Chapter~\ref{ch:background}: it computes FIRST and FOLLOW sets, builds the canonical collection of LR(1) item sets with closure and goto, and fills the ACTION and GOTO tables. For \nova{}'s grammar:

\begin{table}[H]\centering\small
\begin{tabular}{@{}lr@{\hspace{2cm}}lr@{}}\toprule
LR(1) states & \LRStates{} & shift entries & \LRShift{}\\
Productions (excl.\ augmented) & \LRProductions{} & reduce entries & \LRReduce{}\\
Terminals / nonterminals & \LRTerminals{} / \LRNonterminals{} & goto entries & \LRGoto{}\\
\textbf{Conflicts} & \textbf{\LRConflicts{}} & & \\ \bottomrule
\end{tabular}
\caption{Size of the generated parser. The table is built once at start-up from the grammar text, not hard-coded.}
\end{table}

\subsection{A worked example}
Table~\ref{tab:trace} shows the first 26 of \SampleParseSteps{} steps (\SampleParseShifts{} shifts, \SampleParseReduces{} reduces) the parser takes for the program \code{fn main() \{ print(1 + 2 * 3); \}}. The stack alternates \emph{state numbers} and \emph{grammar symbols}; the right-most entries are the top. Watch how \code{1 + 2 * 3} is reduced: the parser shifts \code{2}, \code{*}, \code{3} and reduces \code{MulExpr -> MulExpr * UnaryExpr} \emph{before} reducing the addition, which is precisely how precedence emerges from the table.

\begingroup\small
\begin{longtable}{@{}rp{0.30\textwidth}lp{0.30\textwidth}@{}}\toprule
\textbf{\#} & \textbf{Stack (states and symbols)} & \textbf{Look-ahead} & \textbf{Action}\\ \midrule\endhead
1 & \texttt{\scriptsize 0} & \texttt{fn} & \texttt{\scriptsize shift 4}\\
2 & \texttt{\scriptsize 0 4 fn} & \texttt{ID} & \texttt{\scriptsize shift 6}\\
3 & \texttt{\scriptsize 0 4 fn 6 ID} & \texttt{(} & \texttt{\scriptsize shift 7}\\
4 & \texttt{\scriptsize 0 4 fn 6 ID 7 (} & \texttt{)} & \texttt{\scriptsize reduce 6} \newline {\scriptsize ParamsOpt -\textgreater{} $\varepsilon$}\\
5 & \texttt{\scriptsize 0 4 fn 6 ID 7 ( 8 ParamsOpt} & \texttt{)} & \texttt{\scriptsize shift 12}\\
6 & \texttt{\scriptsize \ldots{} fn 6 ID 7 ( 8 ParamsOpt 12 )} & \texttt{\{} & \texttt{\scriptsize reduce 11} \newline {\scriptsize RetOpt -\textgreater{} $\varepsilon$}\\
7 & \texttt{\scriptsize \ldots{} ID 7 ( 8 ParamsOpt 12 ) 15 RetOpt} & \texttt{\{} & \texttt{\scriptsize shift 25}\\
8 & \texttt{\scriptsize \ldots{} ( 8 ParamsOpt 12 ) 15 RetOpt 25 \{} & \texttt{print} & \texttt{\scriptsize reduce 20} \newline {\scriptsize StmtList -\textgreater{} $\varepsilon$}\\
9 & \texttt{\scriptsize \ldots{} ParamsOpt 12 ) 15 RetOpt 25 \{ 33 StmtList} & \texttt{print} & \texttt{\scriptsize shift 54}\\
10 & \texttt{\scriptsize \ldots{} ) 15 RetOpt 25 \{ 33 StmtList 54 print} & \texttt{(} & \texttt{\scriptsize shift 88}\\
11 & \texttt{\scriptsize \ldots{} RetOpt 25 \{ 33 StmtList 54 print 88 (} & \texttt{INT\_LIT} & \texttt{\scriptsize shift 154}\\
12 & \texttt{\scriptsize \ldots{} \{ 33 StmtList 54 print 88 ( 154 INT\_LIT} & \texttt{+} & \texttt{\scriptsize reduce 69} \newline {\scriptsize Primary -\textgreater{} INT\_LIT}\\
13 & \texttt{\scriptsize \ldots{} \{ 33 StmtList 54 print 88 ( 147 Primary} & \texttt{+} & \texttt{\scriptsize reduce 68} \newline {\scriptsize Postfix -\textgreater{} Primary}\\
14 & \texttt{\scriptsize \ldots{} \{ 33 StmtList 54 print 88 ( 146 Postfix} & \texttt{+} & \texttt{\scriptsize reduce 66} \newline {\scriptsize UnaryExpr -\textgreater{} Postfix}\\
15 & \texttt{\scriptsize \ldots{} \{ 33 StmtList 54 print 88 ( 145 UnaryExpr} & \texttt{+} & \texttt{\scriptsize reduce 63} \newline {\scriptsize MulExpr -\textgreater{} UnaryExpr}\\
16 & \texttt{\scriptsize \ldots{} \{ 33 StmtList 54 print 88 ( 144 MulExpr} & \texttt{+} & \texttt{\scriptsize reduce 59} \newline {\scriptsize AddExpr -\textgreater{} MulExpr}\\
17 & \texttt{\scriptsize \ldots{} \{ 33 StmtList 54 print 88 ( 143 AddExpr} & \texttt{+} & \texttt{\scriptsize shift 267}\\
18 & \texttt{\scriptsize \ldots{} StmtList 54 print 88 ( 143 AddExpr 267 +} & \texttt{INT\_LIT} & \texttt{\scriptsize shift 154}\\
19 & \texttt{\scriptsize \ldots{} print 88 ( 143 AddExpr 267 + 154 INT\_LIT} & \texttt{*} & \texttt{\scriptsize reduce 69} \newline {\scriptsize Primary -\textgreater{} INT\_LIT}\\
20 & \texttt{\scriptsize \ldots{} print 88 ( 143 AddExpr 267 + 147 Primary} & \texttt{*} & \texttt{\scriptsize reduce 68} \newline {\scriptsize Postfix -\textgreater{} Primary}\\
21 & \texttt{\scriptsize \ldots{} print 88 ( 143 AddExpr 267 + 146 Postfix} & \texttt{*} & \texttt{\scriptsize reduce 66} \newline {\scriptsize UnaryExpr -\textgreater{} Postfix}\\
22 & \texttt{\scriptsize \ldots{} print 88 ( 143 AddExpr 267 + 145 UnaryExpr} & \texttt{*} & \texttt{\scriptsize reduce 63} \newline {\scriptsize MulExpr -\textgreater{} UnaryExpr}\\
23 & \texttt{\scriptsize \ldots{} print 88 ( 143 AddExpr 267 + 370 MulExpr} & \texttt{*} & \texttt{\scriptsize shift 269}\\
24 & \texttt{\scriptsize \ldots{} ( 143 AddExpr 267 + 370 MulExpr 269 *} & \texttt{INT\_LIT} & \texttt{\scriptsize shift 154}\\
25 & \texttt{\scriptsize \ldots{} AddExpr 267 + 370 MulExpr 269 * 154 INT\_LIT} & \texttt{)} & \texttt{\scriptsize reduce 69} \newline {\scriptsize Primary -\textgreater{} INT\_LIT}\\
26 & \texttt{\scriptsize \ldots{} AddExpr 267 + 370 MulExpr 269 * 147 Primary} & \texttt{)} & \texttt{\scriptsize reduce 68} \newline {\scriptsize Postfix -\textgreater{} Primary}\\
\bottomrule
\end{longtable}

\endgroup
\captionof{table}{Real parse trace (state numbers are indices into the LR(1) automaton).}\label{tab:trace}

\subsection{Error recovery}\label{sec:recovery}
Stopping at the first syntax error is unfriendly. The grammar therefore contains two \emph{error productions}: \code{Stmt -> error ;} and \code{Stmt -> error}. When the table has no action for the current token the parser:
\begin{enumerate}
  \item reports a diagnostic (suppressing follow-on errors until three tokens have been shifted successfully, to avoid cascades);
  \item pops states until it finds one that can shift the special \code{error} token;
  \item shifts \code{error}, then discards input tokens until one is acceptable in the new state (typically \code{;} or \code{\}});
  \item resumes normal parsing, guaranteeing progress by force-skipping a token if two recoveries happen at the same position.
\end{enumerate}
This is classic \emph{panic mode} (the strategy of yacc/bison). The diagnostics are classified so the message matches the programmer's mistake:

\begin{table}[H]\centering\small
\begin{tabular}{@{}lll@{}}\toprule
\textbf{Code} & \textbf{Situation} & \textbf{Suggested fix}\\ \midrule
E-SYN-001 & next token cannot continue an expression statement & insert \code{;} after the previous token\\
E-SYN-002 & expected \code{)} & insert \code{)}\\
E-SYN-003 & expected \code{\}} or a statement & close the block\\
E-SYN-004 & end of file inside an open construct & add the missing closing token\\
E-SYN-005 & expected an operand & provide a value or variable\\
E-SYN-006 & any other mismatch & list the tokens that \emph{would} be valid here\\ \bottomrule
\end{tabular}
\caption{Syntax diagnostics. The list of expected tokens is read straight from the ACTION table row of the failing state.}
\end{table}

\section{The abstract syntax tree}
The parse tree contains every grammar symbol, including the dozens of unit productions needed to encode precedence. \code{buildAst} converts it to a compact \textbf{typed AST} (\code{Program}, \code{FuncDecl}, \code{Block}, \code{VarDecl}, \code{Assign}, \code{IndexAssign}, \code{If}, \code{While}, \code{Return}, \code{Print}, \code{ExprStmt}, and expression nodes \code{IntLit}, \code{FloatLit}, \code{BoolLit}, \code{StrLit}, \code{VarRef}, \code{Binary}, \code{Unary}, \code{Call}, \code{ArrayLit}, \code{NewArray}, \code{Index}). Every node carries a source range, and a validator checks that children lie inside their parents' ranges and that every expression has a type after analysis.

\section{Semantic analysis}
\subsection{Scopes and symbol tables}
The analyzer walks the AST maintaining a stack of scopes. Each declaration creates a \code{Symbol} (name, type, kind, declaration location, scope id, and lists of read and write locations). The studio's \emph{Semantic Analysis} page displays these tables directly from the analyzer's JSON.

\subsection{Type rules}
\begin{table}[H]\centering\small
\begin{tabular}{@{}lp{0.62\textwidth}@{}}\toprule
\textbf{Construct} & \textbf{Rule}\\ \midrule
\verb|+ - * /| & both operands \code{int} or both \code{float}; result the same type\\
\verb|%| & both \code{int}\\
\verb|< <= > >=| & both numeric of the same type; result \code{bool}\\
\verb|== !=| & both operands the same type (not arrays); result \code{bool}\\
\verb+&&+, \verb+||+, \verb+!+ & \code{bool} operands\\
\code{string + string} & concatenation; other string operators rejected\\
\code{if}, \code{while} condition & \code{bool} (no ``truthiness'' of integers)\\
\code{let x: T = e} & type of \code{e} equals \code{T}\\
call \code{f(a$_1$,\ldots)} & arity matches; each argument's type equals the parameter's type\\
\code{a[i]} & \code{a} is an array, \code{i} is \code{int}; type is the element type\\
\code{print(e)} & \code{e} is not an array\\ \bottomrule
\end{tabular}
\caption{Static typing rules.}
\end{table}

\subsection{Further checks}
\begin{itemize}
  \item \textbf{Missing return}: a function with a result type must return on \emph{every} path; loops are assumed to possibly run zero times.
  \item \textbf{Entry point}: \code{main} must exist with no parameters and no result.
  \item \textbf{Void misuse}: calling a procedure where a value is needed is an error.
  \item \textbf{Warnings}: unused variables (\code{W-SEM-001}) and statements after \code{return} (\code{W-SEM-002}) --- the second never blocks compilation.
\end{itemize}

\subsection{Helpful diagnostics}
Every diagnostic contains a severity, a stable code, a category, the message, an \emph{explanation} of the rule, a \emph{suggestion}, and an exact source range (line, column, end line, end column). For undefined names the suggestion is computed with the \textbf{Levenshtein edit distance}
\[ d(a,b)=\min \begin{cases} d(a_{1..i-1},b)+1 & \text{delete}\\ d(a,b_{1..j-1})+1 & \text{insert}\\ d(a_{1..i-1},b_{1..j-1})+[a_i\ne b_j] & \text{substitute}\end{cases} \]
computed by dynamic programming in $O(|a|\cdot|b|)$; the closest declared name within a length-scaled threshold is offered as ``Did you mean\ldots?''. Section~\ref{sec:quickfix} shows how the assistant turns these suggestions into automatically verified fixes.

\begin{lstlisting}[caption={Compiler output for the shipped example \texttt{e01\_undefined\_variable}.},label={lst:e01},style=tiny]
E-SEM-001  6:11   Variable 'countr' is not declared
                  hint: Did you mean 'counter'?
W-SEM-001  5:9    Variable 'counter' is declared but never read
\end{lstlisting}

\begin{plain}
Notice what the semantic analyzer \emph{does not} do: it never changes the program. It only annotates the tree with types and resolved names, and reports problems. Everything after this point can therefore assume the program is well-typed.
\end{plain}
